% answers_short.tex —— 2017 年·全国Ⅱ卷·理科 · 简短答案速查
% 由 tools/make_short_answers.py 自 answers.tex 的 \gksolhead/\gkanswer 提取，
% 并与 2017/2017_数学真题_全国2理_答案要点.md 逐题比对后生成，勿手改（改源 MD 后重跑）。

\begin{gknotice}
\textbf{使用说明}：本节只给最终答案便于快速核对，完整推导见解析版。答案为编写者独立推导，\textbf{非官方标准答案}。
\end{gknotice}

\gkpart{选择题}
\begin{center}\begin{tabular}{ccccccccccccc}
\toprule
题号 & 1 & 2 & 3 & 4 & 5 & 6 & 7 & 8 & 9 & 10 & 11 & 12 \\
\midrule
答案 & D & C & B & B & A & D & D & B & A & C & A & B \\
\bottomrule
\end{tabular}\end{center}

\gkpart{填空题}
\begin{itemize}[leftmargin=2.6em]
  \item[\textbf{13.}] \(1.96\)（D）
  \item[\textbf{14.}] \(1\)（D）
  \item[\textbf{15.}] \(\dfrac{2n}{n+1}\)（D）
  \item[\textbf{16.}] \(6\)（D）
\end{itemize}

\gkpart{解答题}
\begin{itemize}[leftmargin=2.6em]
  \item[\textbf{17.}] (1) \(\cos B=\dfrac{15}{17}\)（D）；(2) \(b=2\)（D）
  \item[\textbf{18.}] (1) \(P(A)=0.4092\)（D）；(2) 列联表为旧法 \((62,38)\)、新法 \((34,66)\)；\(K^2=\dfrac{1225}{78}\approx15.705>6.635\)，有 99\% 的把握认为产量与方法有关（D）；(3) 新法估计中位数约为 \(52.35\,\mathrm{kg}\)（D）
  \item[\textbf{19.}] (1) \(CE\parallel\) 平面 \(PAB\)（D）；(2) 二面角余弦为 \(\dfrac{\sqrt{10}}5\)（D）
  \item[\textbf{20.}] (1) 点 \(P\) 的轨迹为 \(x^2+y^2=2\)（D）；(2) 过 \(P\) 且垂直于 \(OQ\) 的直线经过椭圆左焦点 \(F=(-1,0)\)（D）
  \item[\textbf{21.}] (1) \(a=1\)（D）；(2) 唯一极大值点 \(x_0\in(0,\tfrac12)\)，且 \(e^{-2}<f(x_0)<\tfrac14\)（D）
  \item[\textbf{22.}] (1) \(x^2+y^2=4x\)，实际轨迹不含原点（即再要求 \(x>0\)）（D）；(2) 三角形面积最大值为 \(2+\sqrt3\)（D）
  \item[\textbf{23.}] (1) \((a+b)(a^5+b^5)\ge4\)（D）；(2) \(a+b\le2\)（D）
\end{itemize}
